% ============================== 3. ПЛАН НА ЕКСПЕРИМЕНТИТЕ ==============================
\chapter{План на експериментите (фаза на машинно обучение и оперативна фаза)}\label{ch:plan}

Експериментите следват работния цикъл на реална IDS и се разделят на две фази: \emph{фаза на машинно обучение}, в която моделите се синтезират, и \emph{оперативна фаза}, в която записаните модели работят върху нови данни (\cref{fig:pipeline}). Около основния експеримент са планирани пет контролни експеримента (E1--E5) девет разширени анализа (A1--A9) и пет проверки за приложимост (A10--A14).

\begin{figure}[!htbp]
\centering
\includegraphics[width=\textwidth]{00_pipeline.png}
\caption{Експериментален конвейер: от суровите CSV файлове до метриките}\label{fig:pipeline}
\vspace{2pt}
\interpret{Горният ред показва подготовката на данните, обща за двете фази. Синята област е фазата на машинно обучение: върху 70\% от данните се обучава един \texttt{Pipeline} (филтър, подбор на признаци, стандартизация, класификатор), който се записва на диска. Оранжевата област е оперативната фаза: записаният модел се зарежда и класифицира останалите 30\%, които не е виждал, на пакети по \num{10000} записа, както би работила IDS върху поток от записи.}
\end{figure}

\section{Фаза на машинно обучение}
За всяка от трите атаки и всеки от трите алгоритъма се изпълняват следните стъпки:
\begin{enumerate}
  \item формиране на балансираната бинарна задача (\cref{eq:binary}) и стратифицирано разделяне 70/30 (\cref{eq:split});
  \item избор на хиперпараметрите с 5-кратна кръстосана валидация и правилото „една стандартна грешка“ (\cref{eq:cv,eq:one-se}). Решетките са: за RF брой дървета $\{200, 300\}$, максимална дълбочина $\{\infty, 25\}$, минимален брой записи в листо $\{2, 5, 10\}$ и дял на признаците при разделяне $\{\sqrt{d}, 0{,}5\}$ (24 комбинации); за XGBoost брой дървета $\{300, 800\}$, дълбочина $\{6, 10\}$, скорост на обучение $\{0{,}03;\ 0{,}1\}$ и \texttt{min\_child\_weight} $\{1, 5\}$ (16 комбинации); за SVM $C\in\{1, 10, 100\}$ и $\gamma\in\{\txt{scale}, 0{,}01;\ 0{,}1\}$ (9 комбинации). RF и XGBoost се настройват върху \num{50000} обучаващи записа, а SVM върху \num{20000}. Всяка комбинация се оценява с целия конвейер;
  \item обучение на финалния конвейер с избраните параметри: RF и XGBoost върху всички обучаващи записи, SVM върху \num{50000} от тях заради сложността $O(n^2)$--$O(n^3)$;
  \item диагностика на преобучението: точност върху обучаващите данни срещу оценката от кръстосаната валидация и точността върху тестовите данни;
  \item запис на конвейера във файл (\texttt{models/<алгоритъм>\_<атака>.joblib}).
\end{enumerate}

\section{Оперативна фаза (фаза на тестване/експлоатация)}
\begin{enumerate}
  \item зареждане на записания модел от диска, без достъп до обучаващите данни;
  \item подаване на тестовите записи (30\%) на последователни пакети по \num{10000} и измерване на времето за класификация;
  \item сравняване на предсказанията с реалните етикети и изчисляване на матрицата на объркване и метриките (\ref{eq:accuracy})--(\ref{eq:auc});
  \item статистическа оценка: bootstrap интервали (\cref{eq:bootstrap}) и тест на McNemar за всяка двойка алгоритми (\cref{eq:mcnemar}).
\end{enumerate}

\section{Контролни експерименти и разширени анализи}
Контролните експерименти проверяват конкретни методологични решения, а разширените анализи обясняват поведението на моделите (\cref{tab:experiments}).

\begin{table}[!htbp]
\centering\small
\caption{План на допълнителните експерименти и анализи}\label{tab:experiments}
\begin{tabularx}{\textwidth}{@{}l p{4.4cm} L@{}}
\toprule
\textsf{\textbf{№}} & \textsf{\textbf{Експеримент}} & \textsf{\textbf{Въпрос, на който отговаря}}\\
\midrule
E1 & Пълен вектор (без редукция) & Каква е цената или ползата от редукцията на размерността?\\
E2 & Тест при естествено съотношение & Колко полезни са алармите при реалния дял на атаките (\cref{eq:baserate})?\\
E3 & Стабилност при кръстосана валидация & Колко се колебае качеството между фолдовете?\\
E4 & Без времевите признаци на потока & Зависят ли моделите от продължителността и интервалите между пакетите, които са най-чувствителни към натоварването на мрежата?\\
E5 & PCA вместо подбор на признаци & Кой метод за редукция е по-подходящ?\\
\midrule
A1 & Профил и KS разделимост & Кои признаци различават атаката от нормалния трафик (\cref{eq:ks})?\\
A2 & 2D проекции (PCA, t-SNE) & Разделими ли са класовете геометрично?\\
A3 & Криви на обучение & Колко данни са нужни и би ли помогнало повече обучение?\\
A4 & Брой запазени измерения & Колко признака или компоненти са необходими и какво печелим по скорост?\\
A5 & Праг при естествено съотношение & Какъв праг осигурява дадено ниво на Precision или Recall?\\
A6 & Подвидове на атаките & Кои разновидности (HOIC, LOIC-UDP, FTP, SSH) се пропускат?\\
A7 & Пренос между атаките & Учат ли моделите конкретната атака или общо „аномалия“?\\
A8 & Пермутационна важност & На кои признаци реално разчита всеки модел?\\
A9 & Bootstrap и McNemar & Значими ли са разликите между алгоритмите?\\
\midrule
A10 & Без признаци-отпечатъци на TCP стека & Разпознават ли моделите атаката или само инструмента и операционната система на атакуващия?\\
A11 & Хронологично разделяне & Открива ли моделът атаката и по-късно във времето, а не само в случайно избрани записи?\\
A12 & Пренос между подвидове & Открива ли модел, обучен върху един инструмент (HOIC, SSH), атака със същата цел, но с друг инструмент?\\
A13 & Нормален трафик от непознат ден & Колко фалшиви аларми дава моделът при нормален трафик, различен от този при обучението?\\
A14 & Комбиниран детектор (XGBoost или SVM) & Може ли съчетанието на двата алгоритъма да открива и нови разновидности на атаката?\\
\bottomrule
\end{tabularx}
\end{table}

Проверките A10--A14 се изпълняват отделно (\texttt{python -m src.robustness}) с хиперпараметрите, избрани в основния експеримент. За тях записите се подреждат по време. В CSV файловете на CSE-CIC-IDS2018 часът е записан в 12-часов формат без означение преди и след обяд. Затова часовете от 1 до 7 се приемат за следобедни: така атаките попадат точно в публикувания график на набора (например LOIC-UDP от 10:09 до 10:43 и HOIC от 14:11 до 14:33 на 21.02.2018)~\cite{sharafaldin2018}. При повтарящ се вектор на признаците се запазва най-ранното му срещане, за да не попадне един и същ запис и в ранните, и в късните данни. Нормалният трафик в A10, A11 и A14 е от деня на съответната атака, за да има смисъл подреждането по време; в A13 се сравнява нормален трафик от един и от два дни.


\paragraph{Детайли на анализите.} При A3 се обучават модели върху $\{10^3, 2\cdot10^3, 5\cdot10^3, 10^4, 2\cdot10^4, 5\cdot10^4, 10^5\}$ записа (SVM до $5\cdot10^4$). При A4 броят запазени измерения е $k\in\{2,5,10,15,20,30,40\}$ за двата метода (подбор на признаци и PCA) плюс контрол без редукция. Двата анализа използват обучаваща извадка до \num{20000} записа (за A4) и обща тестова извадка от \num{10000} невиждани записа. При E2 се използват всички записи, които не са участвали в обучението, смесени в естественото съотношение на класовете в изтеглените файлове (\val{nat-share-DDoS} DDoS, \val{nat-share-BruteForce} BruteForce и \val{nat-share-Bot} Bot). При A5 за всяко целево ниво на Precision се търси прагът $\theta$ с максимален Recall, а за целеви Recall от 0,90 и 0,95 се търси най-високият праг, който го постига. Пермутационната важност (A8) се изчислява върху \num{2000} тестови записа с 3 повторения~\cite{fisher2019}.
